% Lösungen Einheit 1 – Nr. 10–21
\einheitenkopf[le1][1 Kreisumfang]{Einheit 1 von 3 · Kreisumfang}
\erg{10}{a) Radius \quad b) Durchmesser \quad c) keins von beiden (die Strecke geht nicht durch $M$) \quad d) Durchmesser \quad e) Radius}
\erg{11}{a) Umfang \quad b) Fläche \quad c) Fläche \quad d) Umfang \quad e) Umfang}
\erg{12}{a) 6 cm \quad b) 12 m \quad c) 7 cm \quad d) 5 cm \quad e) 7 m \quad f) 4,5 cm \quad g) 4,8 dm \quad h) 0,6 m = 60 cm}
\erg{13}{a) Strecke von $M$ bis zum Rand \quad b) Strecke von Rand zu Rand durch $M$ \quad c) im Mittelpunkt $M$ – jeder Durchmesser geht durch $M$}
\erg{14}{a) Zirkel 2 cm weit öffnen \quad b) $r = 1{,}5\,\text{cm}$ \quad c) $r = 2{,}3\,\text{cm}$}
\erg{15}{Dose 3,14 \quad Teller 3,15 \quad Münze 3,13 \quad Eimer 3,13 \quad a) etwa $3{,}14 \cdot 50\,\text{cm} \approx 157\,\text{cm}$}
\erg{16}{a) 9,4 cm \quad b) 18,8 cm \quad c) 28,3 cm \quad d) 62,8 m \quad e) 25,1 cm \quad f) 11,3 m}
\erg{17}{a) 15,0 cm \quad b) 28,0 m \quad c) $r = 37{,}7 : (2 \cdot \pi) \approx 6{,}0\,\text{cm}$ \quad d) 15,9 cm}
\erg{18}{a) 18,8 cm \quad b) 22,0 m \quad c) $18{,}85 + 12 \approx 30{,}8\,\text{cm}$ \quad d) $u = 2 \cdot \pi \cdot 3{,}8 \approx 23{,}9\,\text{cm}$}
\erg{19}{a) Radius statt Durchmesser eingesetzt; richtig: $u = \pi \cdot 14\,\text{cm} \approx 44{,}0\,\text{cm}$ \quad b) Formel für die Fläche statt für den Umfang; richtig: $u = 2 \cdot \pi \cdot 6\,\text{cm} \approx 37{,}7\,\text{cm}$}
\erg{20}{a) Der Umfang verdoppelt sich, weil $u = \pi \cdot d$: doppelter Durchmesser, doppelter Umfang. \quad b) Nein: Jeder Kreis ist eine vergrößerte oder verkleinerte Kopie jedes anderen; Umfang und Durchmesser wachsen im selben Verhältnis, $u : d$ ist immer $\pi \approx 3{,}14$.}
\erg{21}{a) $\pi \cdot 70\,\text{cm} \approx 219{,}9\,\text{cm}$ \quad b) $500 \cdot 2{,}199\,\text{m} \approx 1099{,}6\,\text{m}$ (gut 1,1 km) \quad c) $5000 : 2{,}199 \approx 2273{,}6$, also etwa 2274 Umdrehungen}
